This section records the present limits of the argument and identifies the exact burden of proof that remains.
It is written for mathematical audit: each unresolved bridge is stated explicitly so that reviewers can test the framework at its most vulnerable points.

\subsection*{Scope and evidentiary standard}
The manuscript proposes a proof framework.
Its intuitive geometric language is intended to guide construction and suggest mechanisms, but it does not by itself constitute proof.
Only statements supported by complete definitions, hypotheses, and logically valid derivations should be treated as formal results.

\subsection*{Primary open bridges}
\begin{enumerate}
  \item \textbf{Rigorous operator construction.}
  The operator/spectral construction must be specified on a precise functional-analytic domain, with well-defined action, closure properties, and spectral meaning.
  At present, this construction is programmatic rather than complete.
  Any downstream claim that depends on the spectrum is therefore conditional on this step.

  \item \textbf{Proof of the zero--balance equivalence.}
  Any claimed equivalence between the proposed balance criteria and nontrivial zeros of $\zeta(s)$ must be proved as a theorem in both directions.
  This equivalence cannot be assumed from analogy, numerics, or partial correspondences.
  Until a full equivalence proof is written, the balance criteria should be treated as candidate reformulations.

  \item \textbf{Rigorous uniqueness of the critical line.}
  The argument requires a proof that the relevant mechanism selects $\Re(s)=\tfrac12$ uniquely, not merely that it is compatible with that line.
  Exclusion of off-line alternatives must be established by explicit estimates or structural theorems under stated hypotheses.
  Without such a proof, uniqueness of the critical line remains open within this framework.
\end{enumerate}

\subsection*{Interpretation of intermediate evidence}
Computational checks, model examples, and heuristic consistency tests are useful for stress-testing definitions and detecting contradictions, but they are not substitutes for proof.
They may justify further investigation; they do not discharge formal obligations.

\subsection*{Policy for unresolved steps}
Any unresolved dependency should be treated as open and cited where it is used.
No conclusion stronger than the proved dependencies is warranted.
Accordingly, the manuscript should be read as materials for mathematical audit of a proposed proof framework, with explicit open bridges rather than implicit assumptions.
